\chapter{The It\^o Integral}

\begin{center}\itshape
A Brownian path has no derivative, so $\int f(B_s)\dd B_s$ makes no sense path by path.\\
It\^o made it make sense on average, and found that the chain rule acquires a second-order term:\\
$\mathrm df(B)=f'(B)\dd B+\frac12f''(B)\dd t$. The roughness of the path is a law of calculus.
\end{center}

In 1942, in a mimeographed journal circulated among Japanese mathematicians, a young statistician of the Cabinet Statistics Bureau in Tokyo, Kiyosi It\^o, proposed to describe a Markov
process not by its transition probabilities, as Kolmogorov and Feller did, but path by path, as the solution of a differential equation driven by Brownian motion. To do so he needed an
integral against $\dd B$, and he built it in a note of 1944, \emph{Stochastic integral}, written while the war made contact with the rest of the world almost impossible. His idea was to
evaluate the integrand at the \emph{left} end of each small interval. With that convention the Riemann sums $\sum f(B_{t_{i-1}})(B_{t_i}-B_{t_{i-1}})$ are martingales, and their limit
exists in $L^2$ although it exists for no single path. The chain rule for this integral, \emph{It\^o's formula}, appeared in 1951. In the foreword to his selected papers (1987), It\^o explained
that he had wanted to look at the paths of Markov processes the way L\'evy looked at sums of independent variables.

Another man had been there first. In February 1940, Wolfgang Doeblin, a twenty-five-year-old French mathematician of German birth serving as a telephone operator in the Ardennes, sent the
Acad\'emie des Sciences a sealed envelope, a \emph{pli cachet\'e}, titled \emph{Sur l'\'equation de Kolmogoroff}. He died in June 1940, by his own hand, as his unit was about to be captured.
The envelope stayed sealed for sixty years. When it was opened in 2000, at the request of his brother, it was found to contain a construction of diffusions by a random time change of Brownian
motion and a version of It\^o's formula. Bernard Bru and Marc Yor told the story in 2002.

This chapter builds the integral $\int H\dd M$ against any continuous local martingale $M$, proves It\^o's formula in several variables (the formula promised in Volume~4, Chapter~11), and draws
the first consequences, each of which became a theorem with a name: L\'evy's characterisation of Brownian motion by its clock, the theorem of Dambis, Dubins and Schwarz that every continuous
martingale is a time-changed Brownian motion, the inequalities of Burkholder, and Tanaka's formula for $|B_t|$, which introduces the \emph{local time} at $0$ and proves the theorem of L\'evy announced
in Chapter~2: $M-B$ and $|B|$ have the same law as processes. With it, the zero set of Brownian motion has dimension exactly one half.

\begin{reminder}[from Chapter 3]
Filtrations satisfy the usual conditions. For continuous local martingales $M,N$ with $M_0=N_0=0$: the bracket $\langle M,N\rangle$ is the unique continuous adapted finite variation process with
$MN-\langle M,N\rangle$ a local martingale, and $\sum\Delta M\Delta N\to\langle M,N\rangle$ in probability along partitions (Propositions~\ref{prop:qvsums} and~\ref{prop:bracket}). A continuous local martingale of finite
variation is constant (Theorem~\ref{thm:fvconstant}). If $\mathbb E\langle M\rangle_\infty<\infty$, $M$ is a martingale bounded in $L^2$ with $\mathbb EM_\infty^2=\mathbb E\langle M\rangle_\infty$, and conversely
(Corollary~\ref{cor:L2bracket}). A non-negative continuous local martingale with integrable initial value is a supermartingale (Exercise~\ref{ex:superm}). For a bounded martingale and partitions with
mesh tending to $0$, $\sum_i\mathbb E\bigl[\bigl((\Delta M_i)^2-\Delta\langle M\rangle_i\bigr)^2\bigr]\to0$, and these terms are orthogonal (proof of Proposition~\ref{prop:qvsums}).
\end{reminder}

\section{Brackets, stopping and the Kunita--Watanabe inequality}

Two technical facts about brackets come first.

\begin{lemma}\label{lem:bracketstop}
Let $M,N$ be continuous local martingales with $M_0=N_0=0$ and $\tau$ a stopping time. Then $\langle M^\tau,N\rangle=\langle M,N\rangle^\tau$.
\end{lemma}

\begin{proof}
By stopping $M$ and $N$ at $\inf\{t:|M_t|+|N_t|\ge n\}$ we may assume them bounded (Proposition~\ref{prop:localmart}). Write $M^\tau N-\langle M,N\rangle^\tau=(MN-\langle M,N\rangle)^\tau+Y$ with
$Y_t=M_{\tau\wedge t}(N_t-N_{\tau\wedge t})$. The first term is a martingale (Corollary~\ref{cor:stopped}). For $Y$, let $s<t$ and $A\in\mathcal F_s$. On $C=A\cap\{\tau\le s\}$, $Y_t-Y_s=M_{\tau\wedge s}(N_t-N_s)$ and
$M_{\tau\wedge s}\mathbf 1_C$ is $\mathcal F_s$-measurable, so $\mathbb E[(Y_t-Y_s)\mathbf 1_C]=0$. On $D=A\cap\{\tau>s\}$, $Y_s=0$; as shown in the proof of Corollary~\ref{cor:stopped}, $D\in\mathcal F_{\tau\wedge s}\subseteq\mathcal F_{\tau\wedge t}$,
so $\mathbb E[Y_t\mathbf 1_D]=\mathbb E\bigl[M_{\tau\wedge t}\mathbf 1_D\,\mathbb E[N_t-N_{\tau\wedge t}\mid\mathcal F_{\tau\wedge t}]\bigr]=0$ by optional stopping (Theorem~\ref{thm:ostcont}). So $Y$ is a martingale, and uniqueness of
the bracket concludes.
\end{proof}

\begin{lemma}[Kunita--Watanabe inequality, 1967]\label{lem:kw}
Let $M,N$ be continuous local martingales with $M_0=N_0=0$, and $H,K$ measurable processes. Almost surely,
\[
\int_0^\infty|H_sK_s|\,|\mathrm d\langle M,N\rangle_s|\le\Bigl(\int_0^\infty H_s^2\dd\langle M\rangle_s\Bigr)^{1/2}\Bigl(\int_0^\infty K_s^2\dd\langle N\rangle_s\Bigr)^{1/2},
\]
where $|\mathrm d\langle M,N\rangle|$ is the total variation measure of $\langle M,N\rangle$.
\end{lemma}

\begin{proof}
For rationals $s<t$ and $r$, $\langle M+rN\rangle_t-\langle M+rN\rangle_s\ge0$ almost surely, that is $a+2rc+r^2b\ge0$ with $a,b,c$ the increments of $\langle M\rangle$, $\langle N\rangle$, $\langle M,N\rangle$ over $[s,t]$.
Outside a null set this holds for all rationals, hence (by continuity) for all real $r$ and all $s<t$, and the discriminant gives $|c|\le\sqrt{ab}$. Fix $\omega$ outside this null set and $T>0$, and on $[0,T]$ let
$\alpha=\mathrm d\langle M\rangle$, $\nu=\mathrm d\langle N\rangle$ and $\mu=|\mathrm d\langle M,N\rangle|$. For an interval $I$, splitting it into pieces $I_k$ and using Cauchy--Schwarz,
$\sum_k|c(I_k)|\le\sum_k\sqrt{\alpha(I_k)\nu(I_k)}\le\sqrt{\alpha(I)\nu(I)}$, so $\mu(I)\le\sqrt{\alpha(I)\nu(I)}\le\frac12\bigl(r^2\alpha(I)+r^{-2}\nu(I)\bigr)$ for every $r>0$.

This inequality between the measures $\mu$ and $\frac12(r^2\alpha+r^{-2}\nu)$ extends from intervals to open sets (countable disjoint unions of intervals) and then to all Borel sets $A$, by outer regularity of finite
Borel measures (Volume~7, Lemma~11.10). Minimising over $r$ gives $\mu(A)\le\sqrt{\alpha(A)\nu(A)}$. For simple functions $f=\sum_jc_j\mathbf 1_{A_j}$ and $g=\sum_jd_j\mathbf 1_{A_j}$ on a common partition,
\[
\int|fg|\dd\mu\le\sum_j|c_jd_j|\sqrt{\alpha(A_j)\nu(A_j)}\le\Bigl(\sum_jc_j^2\alpha(A_j)\Bigr)^{1/2}\Bigl(\sum_jd_j^2\nu(A_j)\Bigr)^{1/2}.
\]
Monotone limits give the inequality for $|H(\cdot,\omega)|$ and $|K(\cdot,\omega)|$ on $[0,T]$; let $T\to\infty$.
\end{proof}

\section{The integral}

\begin{definition}\label{def:H2}
$\mathbb H^2$ is the space of continuous martingales $M$ with $M_0=0$ that are bounded in $L^2$, with the norm $\|M\|_{\mathbb H^2}=(\mathbb EM_\infty^2)^{1/2}=(\mathbb E\langle M\rangle_\infty)^{1/2}$. For $M\in\mathbb H^2$,
$L^2(M)$ is the space of progressive processes $H$ with $\|H\|^2_{L^2(M)}=\mathbb E\int_0^\infty H_s^2\dd\langle M\rangle_s<\infty$ (up to equality $\mathrm d\langle M\rangle\,\mathrm dP$-almost everywhere). An \emph{elementary process}
is $H=\sum_{i=0}^{p-1}H_i\mathbf 1_{]t_i,t_{i+1}]}$ with $0\le t_0<\dots<t_p$ and $H_i$ bounded and $\mathcal F_{t_i}$-measurable.
\end{definition}

The two norms on $\mathbb H^2$ agree by Corollary~\ref{cor:L2bracket}. A progressive process is in particular measurable on $[0,\infty[\times\Omega$, and $\int_0^tH_s\dd A_s$ is then $\mathcal F_t$-measurable for every
continuous adapted increasing $A$ (a monotone class argument, starting from $H=\mathbf 1_{[0,u]\times C}$, $C\in\mathcal F_t$).

\begin{lemma}\label{lem:H2complete}
$\mathbb H^2$ is a Hilbert space, and $L^2(M)$ is a Hilbert space in which elementary processes are dense.
\end{lemma}

\begin{proof}
Let $(M^n)$ be Cauchy in $\mathbb H^2$. Then $M^n_\infty\to Y$ in $L^2$, and by Doob's inequality $\mathbb E\sup_t(M^n_t-M^m_t)^2\le4\|M^n-M^m\|^2$, so a subsequence converges uniformly almost surely to a continuous
adapted $M$, with $M_t=\lim M^n_t=\mathbb E[Y\mid\mathcal F_t]$ in $L^2$. So $M\in\mathbb H^2$, $M_\infty=Y$ (Theorem~\ref{thm:L2cont}) and $M^n\to M$. The space $L^2(M)$ is an $L^2$ space for the measure
$\mathrm d\langle M\rangle\,\mathrm dP$ restricted to the progressive $\sigma$-algebra, hence complete (Volume~7, Theorem~10.4).

For density, let $K\in L^2(M)$ be orthogonal to every elementary process. Since $\mathbb E\int|K|\dd\langle M\rangle\le\|K\|(\mathbb E\langle M\rangle_\infty)^{1/2}<\infty$, the process $X_t=\int_0^tK_s\dd\langle M\rangle_s$ is
continuous, adapted, integrable and of finite variation. Orthogonality to $F\mathbf 1_{]s,t]}$, with $F$ bounded and $\mathcal F_s$-measurable, says $\mathbb E[F(X_t-X_s)]=0$: $X$ is a martingale. By
Theorem~\ref{thm:fvconstant}, $X=0$, so the signed measure $K\dd\langle M\rangle$ vanishes on intervals, hence everywhere (Volume~7, Theorem~7.5, applied to its positive and negative parts), and
$\|K\|^2=\mathbb E\int K\cdot K\dd\langle M\rangle=0$.
\end{proof}

For an elementary $H$ and $M\in\mathbb H^2$ set $(H\cdot M)_t=\sum_iH_i\bigl(M_{t_{i+1}\wedge t}-M_{t_i\wedge t}\bigr)$, and for a process $A$ of finite variation let $(H\cdot A)_t=\int_0^tH_s\dd A_s$.

\begin{lemma}\label{lem:elementary}
For elementary $H$ and $M,N\in\mathbb H^2$: $H\cdot M\in\mathbb H^2$, $\langle H\cdot M,N\rangle=H\cdot\langle M,N\rangle$, and $\|H\cdot M\|_{\mathbb H^2}=\|H\|_{L^2(M)}$.
\end{lemma}

\begin{proof}
Fix $i$ and let $Z=(M^{t_{i+1}}-M^{t_i})N-\bigl(\langle M,N\rangle^{t_{i+1}}-\langle M,N\rangle^{t_i}\bigr)$, a martingale by Lemma~\ref{lem:bracketstop}, which vanishes on $[0,t_i]$. Then $H_iZ$ is a martingale: for
$s\ge t_i$ take $H_i$ out of the conditional expectation, and for $s<t_i$ condition first on $\mathcal F_{t_i}$. Summing over $i$, $(H\cdot M)N-H\cdot\langle M,N\rangle$ is a martingale. Similarly $H\cdot M$ is a
martingale. With $N=M$ this gives $\langle H\cdot M,M\rangle=H\cdot\langle M\rangle$; with $N=H\cdot M$ and symmetry, $\langle H\cdot M\rangle=H\cdot\langle M,H\cdot M\rangle=H\cdot(H\cdot\langle M\rangle)=H^2\cdot\langle M\rangle$.
This is bounded by $(\max_i\|H_i\|_\infty)^2\langle M\rangle_\infty$, which is integrable, so Corollary~\ref{cor:L2bracket} gives $H\cdot M\in\mathbb H^2$ and $\|H\cdot M\|^2=\mathbb E\langle H\cdot M\rangle_\infty=\|H\|^2_{L^2(M)}$.
\end{proof}

\begin{theorem}[the It\^o integral]\label{thm:itoH2}
Let $M\in\mathbb H^2$ and $H\in L^2(M)$. There is a unique $H\cdot M\in\mathbb H^2$ such that
\[
\langle H\cdot M,N\rangle=H\cdot\langle M,N\rangle\qquad\text{for every }N\in\mathbb H^2 .
\]
The map $H\mapsto H\cdot M$ is a linear isometry from $L^2(M)$ to $\mathbb H^2$ that extends the elementary integral. We write $(H\cdot M)_t=\int_0^tH_s\dd M_s$.
\end{theorem}

\begin{proof}
\emph{Uniqueness.} If $X,X'$ qualify, $\langle X-X',N\rangle=0$ for all $N$; with $N=X-X'$, $\mathbb E(X_\infty-X'_\infty)^2=\mathbb E\langle X-X'\rangle_\infty=0$, and Doob's inequality gives $X=X'$.

\emph{Existence.} By Lemma~\ref{lem:H2complete} choose elementary $H^n\to H$ in $L^2(M)$. By Lemma~\ref{lem:elementary}, $(H^n\cdot M)$ is Cauchy in $\mathbb H^2$; let $X$ be its limit. Fix $N\in\mathbb H^2$ and $t$. By
Lemma~\ref{lem:kw} and Cauchy--Schwarz,
\[
\mathbb E\bigl|\langle H^n\cdot M-X,N\rangle_t\bigr|\le\|H^n\cdot M-X\|\,\|N\|,\qquad\mathbb E\Bigl|\int_0^t(H^n_s-H_s)\dd\langle M,N\rangle_s\Bigr|\le\|H^n-H\|\,\|N\|,
\]
and both tend to $0$. Since $\langle H^n\cdot M,N\rangle_t=\int_0^tH^n\dd\langle M,N\rangle$, we get $\langle X,N\rangle_t=\int_0^tH\dd\langle M,N\rangle$ almost surely for each $t$, hence for all $t$ by continuity. The
isometry passes to the limit, and linearity follows from uniqueness.
\end{proof}

\begin{proposition}\label{prop:itoprops}
Let $M,N\in\mathbb H^2$, $H\in L^2(M)$, $K\in L^2(N)$ and $\tau$ a stopping time.
\begin{enumerate}[label=(\alph*)]
\item $\langle H\cdot M,K\cdot N\rangle=HK\cdot\langle M,N\rangle$. In particular $\langle H\cdot M\rangle=H^2\cdot\langle M\rangle$ and $\mathbb E\bigl[(\int_0^tH\dd M)^2\bigr]=\mathbb E\int_0^tH_s^2\dd\langle M\rangle_s$.
\item If $G$ is progressive and $GH\in L^2(M)$, then $G\in L^2(H\cdot M)$ and $G\cdot(H\cdot M)=GH\cdot M$.
\item $(H\cdot M)^\tau=H\mathbf 1_{[0,\tau]}\cdot M=H\cdot M^\tau$.
\end{enumerate}
\end{proposition}

\begin{proof}
All three follow from the uniqueness in Theorem~\ref{thm:itoH2} and the associativity of Stieltjes integrals, $G\cdot(H\cdot A)=GH\cdot A$. (a) $\langle H\cdot M,K\cdot N\rangle=H\cdot\langle M,K\cdot N\rangle=H\cdot(K\cdot\langle M,N\rangle)$.
(b) $\int G^2\dd\langle H\cdot M\rangle=\int G^2H^2\dd\langle M\rangle$, and $\langle G\cdot(H\cdot M),N\rangle=G\cdot(H\cdot\langle M,N\rangle)=GH\cdot\langle M,N\rangle$. (c) By Lemma~\ref{lem:bracketstop},
$\langle(H\cdot M)^\tau,N\rangle=(H\cdot\langle M,N\rangle)^\tau=H\mathbf 1_{[0,\tau]}\cdot\langle M,N\rangle$, and also $=H\cdot\langle M,N\rangle^\tau=H\cdot\langle M^\tau,N\rangle$.
\end{proof}

\section{Local martingales, semimartingales and Riemann sums}

\begin{theorem}\label{thm:itoloc}
Let $M$ be a continuous local martingale with $M_0=0$, and $H$ a progressive process with $\int_0^tH_s^2\dd\langle M\rangle_s<\infty$ for all $t$, almost surely. There is a unique continuous local martingale
$H\cdot M$, zero at $0$, with $\langle H\cdot M,N\rangle=H\cdot\langle M,N\rangle$ for every continuous local martingale $N$. Proposition~\ref{prop:itoprops} holds in this setting. If
$\mathbb E\int_0^tH^2\dd\langle M\rangle<\infty$, then $(H\cdot M)^t\in\mathbb H^2$ and $\mathbb E(\int_0^tH\dd M)^2=\mathbb E\int_0^tH^2\dd\langle M\rangle$.
\end{theorem}

\begin{proof}
Let $\tau_n=\inf\{t:\int_0^t(1+H_s^2)\dd\langle M\rangle_s\ge n\}$, which increases to $\infty$. Then $\langle M^{\tau_n}\rangle\le n$, so $M^{\tau_n}\in\mathbb H^2$, and $H\in L^2(M^{\tau_n})$. Let $X^n=H\cdot M^{\tau_n}$. By
Proposition~\ref{prop:itoprops}(c), $(X^{n+1})^{\tau_n}=H\cdot M^{\tau_n}=X^n$, so there is a continuous process $X$ with $X^{\tau_n}=X^n$ for all $n$: a continuous local martingale. For a continuous local martingale $N$,
stopped at $\sigma_m$ so that $N^{\sigma_m}\in\mathbb H^2$, Lemma~\ref{lem:bracketstop} gives $\langle X,N\rangle^{\tau_n\wedge\sigma_m}=\langle X^n,N^{\sigma_m}\rangle=H\cdot\langle M^{\tau_n},N^{\sigma_m}\rangle=(H\cdot\langle M,N\rangle)^{\tau_n\wedge\sigma_m}$;
let $n,m\to\infty$. If $X'$ also qualifies, $D=X-X'$ has $\langle D\rangle=0$, so $D^2$ is a non-negative local martingale with $D_0=0$, a supermartingale with mean $0$ (Exercise~\ref{ex:superm}): $D=0$. The
properties follow by the same localisation, and the last claim from Corollary~\ref{cor:L2bracket}.
\end{proof}

\begin{definition}\label{def:semimart}
A \emph{continuous semimartingale} is a process $X=X_0+M+A$ with $X_0$ $\mathcal F_0$-measurable, $M$ a continuous local martingale and $A$ a continuous adapted process of finite variation, $M_0=A_0=0$. The
decomposition is unique (Theorem~\ref{thm:fvconstant}). For two semimartingales, $\langle X,Y\rangle=\langle M,N\rangle$. For $H$ progressive and \emph{locally bounded} ($\sup_{s\le t}|H_s|<\infty$ for each $t$, for example
continuous and adapted), $\int_0^tH\dd X=\int_0^tH\dd M+\int_0^tH\dd A$.
\end{definition}

Brownian motion is a semimartingale, and so is $f(t,B_t)$ for smooth $f$, as It\^o's formula will show; the integrals above are the only ones we need. The next result says that It\^o's integral is the limit
of the left-point Riemann sums, which is how It\^o defined it.

\begin{proposition}[Riemann sums]\label{prop:riemann}
Let $X$ be a continuous semimartingale, $H$ a continuous adapted process, and $\pi_n$ partitions of $[0,t]$ with mesh tending to $0$. Then, in probability,
\[
\sum_iH_{t_{i-1}}\bigl(X_{t_i}-X_{t_{i-1}}\bigr)\longrightarrow\int_0^tH\dd X,\qquad\sum_iH_{t_{i-1}}\bigl(X_{t_i}-X_{t_{i-1}}\bigr)^2\longrightarrow\int_0^tH\dd\langle X\rangle .
\]
By polarisation, $\sum_iH_{t_{i-1}}\Delta X_i\Delta Y_i\to\int_0^tH\dd\langle X,Y\rangle$ for two semimartingales.
\end{proposition}

\begin{proof}
Let $\tau_k=\inf\{s:|H_s|+|M_s|+\langle M\rangle_s+V_s\ge k\}$, with $V$ the variation of $A$. All the quantities agree on $\{\tau_k>t\}$ with those of the stopped processes, and $P(\tau_k\le t)\to0$; so we may assume
$|H|,|M|,\langle M\rangle,V\le k$. The process $H^\pi=\sum_iH_{t_{i-1}}\mathbf 1_{]t_{i-1},t_i]}$ is elementary, $\int_0^tH^\pi\dd M$ is the martingale part of the first sum, and $H^\pi\to H$ pointwise with $|H^\pi|\le k$. By the isometry,
$\mathbb E(\int_0^t(H^\pi-H)\dd M)^2=\mathbb E\int_0^t(H^\pi-H)^2\dd\langle M\rangle\to0$ by dominated convergence, and $\int_0^tH^\pi\dd A\to\int_0^tH\dd A$ path by path.

For the second sum, $(\Delta X_i)^2=(\Delta M_i)^2+2\Delta M_i\Delta A_i+(\Delta A_i)^2$, and the last two terms contribute at most $k(2\max_i|\Delta M_i|+\max_i|\Delta A_i|)V_t\to0$. Let $Z_i=(\Delta M_i)^2-\Delta\langle M\rangle_i$.
The terms $H_{t_{i-1}}Z_i$ are orthogonal, since $\mathbb E[Z_i\mid\mathcal F_{t_{i-1}}]=0$, so $\mathbb E(\sum_iH_{t_{i-1}}Z_i)^2\le k^2\sum_i\mathbb EZ_i^2\to0$ (reminder). Finally
$\sum_iH_{t_{i-1}}\Delta\langle M\rangle_i\to\int_0^tH\dd\langle M\rangle$, a Riemann--Stieltjes sum of a continuous function.
\end{proof}

\section{It\^o's formula}

\begin{theorem}[It\^o's formula, 1951]\label{thm:itoformula}
Let $X=(X^1,\dots,X^d)$ be continuous semimartingales and $f\in C^2(\R^d)$. Almost surely, for all $t\ge0$,
\[
f(X_t)=f(X_0)+\sum_{k=1}^d\int_0^t\partial_kf(X_s)\dd X^k_s+\frac12\sum_{k,l=1}^d\int_0^t\partial_k\partial_lf(X_s)\dd\langle X^k,X^l\rangle_s .
\]
In particular $f(X)$ is a continuous semimartingale.
\end{theorem}

\begin{proof}
Both sides are continuous in $t$, so it suffices to prove the identity almost surely for each fixed $t$. Stopping at $\inf\{s:|X_s|\ge k\}$ (and letting $k\to\infty$) we may assume that $X$ stays in a ball $\bar B(0,k)$, on
which $f$ and its first two derivatives are bounded and the second derivatives have a modulus of uniform continuity $\omega$. Taylor's formula with integral remainder gives, for $x,y$ in the ball,
\[
f(y)-f(x)=\sum_k\partial_kf(x)(y-x)_k+\frac12\sum_{k,l}\partial_k\partial_lf(x)(y-x)_k(y-x)_l+R(x,y),
\]
with $|R(x,y)|\le\frac{d^2}2\,\omega(|y-x|)\,|y-x|^2$.
Apply it along a partition of $[0,t]$ with $x=X_{t_{i-1}}$, $y=X_{t_i}$, and sum over $i$. By Proposition~\ref{prop:riemann}, the first-order terms converge in probability to $\sum_k\int_0^t\partial_kf(X)\dd X^k$ and the
second-order terms to $\frac12\sum_{k,l}\int_0^t\partial_k\partial_lf(X)\dd\langle X^k,X^l\rangle$. The remainders sum to at most $\frac{d^2}2\,\omega(\max_i|\Delta X_i|)\sum_i|\Delta X_i|^2$; the sum converges in probability to
$\sum_k\langle X^k\rangle_t$ and $\omega(\max_i|\Delta X_i|)\to0$ by uniform continuity of the paths. So the total tends to $0$ in probability, and the identity holds almost surely.
\end{proof}

For one Brownian motion this is the formula of Volume~4, Chapter~11: $f(B_t)=f(0)+\int_0^tf'(B_s)\dd B_s+\frac12\int_0^tf''(B_s)\dd s$. With time as an extra coordinate ($X^0_t=t$, of finite variation, so all its
brackets vanish), $f(t,B_t)=f(0,0)+\int_0^t\partial_xf\dd B+\int_0^t(\partial_t+\frac12\partial_x^2)f\dd s$: the heat operator appears, and $f(t,B_t)$ is a local martingale exactly when $f$ solves the backward heat equation.
Complex-valued $f$ are handled by taking real and imaginary parts.

\begin{think}
Compute $\int_0^tB_s\dd B_s$ with $f(x)=x^2$. Ordinary calculus would give $\frac12B_t^2$. Which sums converge to that answer instead (Exercise~\ref{ex:bdb})?
\end{think}

\begin{figure}[ht]
\centering
\fig{ch04-sums}
\caption{Along one Brownian path on $[0,2]$ with $2^{12}$ steps, the running sums $\sum B_{t_{i-1}}\Delta B_i$ (It\^o, left point), $\sum\frac12(B_{t_{i-1}}+B_{t_i})\Delta B_i$ (Stratonovich, midpoint) and $\sum B_{t_i}\Delta B_i$
(right point), against $\frac12(B_t^2-t)$, $\frac12B_t^2$ and $\frac12(B_t^2+t)$ (thin black). The three conventions differ by $\frac12\sum(\Delta B_i)^2\approx\frac t2$: the quadratic variation.}
\label{fig:sums}
\end{figure}

\section{L\'evy's characterisation and time changes}

\begin{theorem}[L\'evy's characterisation, 1948]\label{thm:levychar}
Let $M=(M^1,\dots,M^d)$ be continuous local martingales with $M_0=0$ and $\langle M^k,M^l\rangle_t=\delta_{kl}t$. Then $M$ is a $d$-dimensional $(\mathcal F_t)$-Brownian motion.
\end{theorem}

\begin{proof}
Fix $\xi\in\R^d$ and let $f(t,x)=\exp\bigl(i\xi\cdot x+\frac12|\xi|^2t\bigr)$. It\^o's formula, with $\partial_tf=\frac12|\xi|^2f$ and $\frac12\sum_k\partial_k^2f=-\frac12|\xi|^2f$, gives
$f(t,M_t)=1+\sum_ki\xi_k\int_0^tf(s,M_s)\dd M^k_s$, a complex continuous local martingale. It is bounded by $e^{|\xi|^2T/2}$ on $[0,T]$, hence a martingale (Proposition~\ref{prop:localmart}(c)), and for $s<t$
\[
\mathbb E\bigl[e^{i\xi\cdot(M_t-M_s)}\mid \mathcal F_s\bigr]=e^{-\frac12|\xi|^2(t-s)} .
\]
So for $A\in\mathcal F_s$ with $P(A)>0$, the law of $M_t-M_s$ under $P(\cdot\mid A)$ has the characteristic function of $N(0,(t-s)I_d)$: the increment is Gaussian with this covariance and independent of $\mathcal F_s$.
\end{proof}

The theorem is remarkable for what it does not assume: no Gaussian law, no independence, only continuity and the clock. A fair game that moves continuously and whose squared fluctuations accumulate at
unit rate \emph{is} Brownian motion. The same computation with a real exponent gives the exponential martingales of Chapter~3 for every local martingale.

\begin{proposition}[exponential martingales]\label{prop:expmartgen}
Let $M$ be a continuous local martingale with $M_0=0$ and $\lambda\in\R$. Then $\mathcal E^\lambda(M)_t=\exp\bigl(\lambda M_t-\frac{\lambda^2}2\langle M\rangle_t\bigr)=1+\lambda\int_0^t\mathcal E^\lambda(M)_s\dd M_s$ is a continuous local martingale and
a supermartingale. For $a,c>0$,
\[
P\Bigl(\sup_{t\ge0}M_t\ge a,\ \langle M\rangle_\infty\le c\Bigr)\le e^{-a^2/2c}.
\]
\end{proposition}

\begin{proof}
It\^o's formula for $f(x,y)=e^{\lambda x-\lambda^2y/2}$ with $X=(M,\langle M\rangle)$ gives $\mathrm df=\lambda f\dd M-\frac{\lambda^2}2f\dd\langle M\rangle+\frac{\lambda^2}2f\dd\langle M\rangle$; it is a supermartingale by
Exercise~\ref{ex:superm}. Let $\tau=\inf\{t:M_t\ge a\}$. The stopped process $\mathcal E^\lambda(M)^\tau$ is again a non-negative local martingale, so $\mathbb E\,\mathcal E^\lambda(M)_{\tau\wedge t}\le1$. On the event
$\{\tau\le t,\ \langle M\rangle_\infty\le c\}$, $\mathcal E^\lambda(M)_{\tau\wedge t}\ge e^{\lambda a-\lambda^2c/2}$. Hence $P(\tau\le t,\langle M\rangle_\infty\le c)\le e^{-\lambda a+\lambda^2c/2}$. Let $t\to\infty$ and take $\lambda=a/c$.
\end{proof}

This is the continuous version of the Azuma--Hoeffding inequality of Volume~4, Theorem~10.24, and it is sharp: for $M=B^{c}$ (Brownian motion stopped at time $c$) the left side is $P(\max_{s\le c}B_s\ge a)=2P(B_c\ge a)$,
which is $e^{-a^2/2c}$ up to a polynomial factor.

The clock of Chapter~3 can be used to change time. Two lemmas prepare the theorem.

\begin{lemma}\label{lem:ostL2}
Let $X$ be a continuous martingale with $\mathbb E\sup_t|X_t|<\infty$ which converges almost surely to a limit $X_\infty$ (for instance, a martingale bounded in $L^2$). If $\sigma\le\tau$ are stopping times, possibly infinite, then $\mathbb E[X_\tau\mid\mathcal F_\sigma]=X_\sigma$.
\end{lemma}

\begin{proof}
For $A\in\mathcal F_\sigma$, $A\cap\{\sigma\le n\}\in\mathcal F_{\sigma\wedge n}$ (check on $\{\sigma\wedge n\le u\}$), so Theorem~\ref{thm:ostcont} gives $\mathbb E[X_{\tau\wedge n}\mathbf 1_{A\cap\{\sigma\le n\}}]=\mathbb E[X_{\sigma\wedge n}\mathbf 1_{A\cap\{\sigma\le n\}}]$. All
variables are dominated by the integrable $\sup_t|X_t|$. Letting $n\to\infty$, $\mathbb E[X_\tau\mathbf 1_{A\cap\{\sigma<\infty\}}]=\mathbb E[X_\sigma\mathbf 1_{A\cap\{\sigma<\infty\}}]$. On $\{\sigma=\infty\}$, $X_\tau=X_\infty=X_\sigma$.
\end{proof}

\begin{lemma}\label{lem:constancy}
Let $M$ be a continuous local martingale. Almost surely, $M$ is constant on every interval on which $\langle M\rangle$ is constant.
\end{lemma}

\begin{proof}
For rational $q$ let $\sigma_q=\inf\{t>q:\langle M\rangle_t>\langle M\rangle_q\}$ and $N=\mathbf 1_{]q,\sigma_q]}\cdot M=M^{\sigma_q}-M^q$. Then $\langle N\rangle=\mathbf 1_{]q,\sigma_q]}\cdot\langle M\rangle=0$, so $N=0$ (as in the proof of
Theorem~\ref{thm:itoloc}): $M$ is constant on $[q,\sigma_q]$, for all rational $q$ at once, almost surely. If $\langle M\rangle$ is constant on $[s,t]$ with $s<t$, then $\sigma_q\ge t$ for every rational $q\in\,]s,t[$, so $M$ is
constant on $[q,t]$; let $q\downarrow s$.
\end{proof}

\begin{theorem}[Dambis, Dubins--Schwarz, 1965]\label{thm:dds}
Let $M$ be a continuous local martingale with $M_0=0$ and $\langle M\rangle_\infty=\infty$ almost surely. Let $T_s=\inf\{t:\langle M\rangle_t>s\}$. Then $\beta_s=M_{T_s}$ is a Brownian motion for the filtration
$(\mathcal F_{T_s})_{s\ge0}$, and $M_t=\beta_{\langle M\rangle_t}$ for all $t$.
\end{theorem}

\begin{proof}
Each $T_s$ is a finite stopping time ($\{T_s<t\}=\{\langle M\rangle_t>s\}$, and the filtration is right-continuous), $s\mapsto T_s$ is increasing and right-continuous, and $\langle M\rangle_{T_s}=s$. The filtration
$\mathcal G_s=\mathcal F_{T_s}$ satisfies the usual conditions: if $A\in\bigcap_n\mathcal G_{s+1/n}$, then $A\cap\{T_s<t\}=\bigcup_nA\cap\{T_{s+1/n}<t\}\in\mathcal F_t$, and so $A\cap\{T_s\le t\}\in\mathcal F_{t+}=\mathcal F_t$.

\emph{Continuity.} If $T_{s-}<T_s$, then $\langle M\rangle=s$ on $[T_{s-},T_s]$, so $M$ is constant there by Lemma~\ref{lem:constancy}, and $\beta_{s-}=M_{T_{s-}}=M_{T_s}=\beta_s$. So $\beta$ is continuous, and adapted to $(\mathcal G_s)$
by Lemma~\ref{lem:stopping}(b).

\emph{Brownian motion.} Fix $u>0$. The local martingale $M^{T_u}$ has bracket $\langle M\rangle^{T_u}\le u$, so by Corollary~\ref{cor:L2bracket} it is bounded in $L^2$ and $(M^{T_u})^2-\langle M\rangle^{T_u}$ is a uniformly integrable
martingale. Both converge as $t\to\infty$ (Theorem~\ref{thm:L2cont}, and $\langle M\rangle^{T_u}$ is increasing and bounded) and are dominated by $\sup_tM_{T_u\wedge t}^2+u$, which is integrable by Doob's inequality. Lemma~\ref{lem:ostL2}, applied to $T_s\le T_t$ for $s<t\le u$, gives
$\mathbb E[\beta_t\mid\mathcal G_s]=\beta_s$ and $\mathbb E[\beta_t^2-t\mid\mathcal G_s]=\beta_s^2-s$. So $\beta$ is a continuous $(\mathcal G_s)$-martingale with $\langle\beta\rangle_s=s$: a Brownian motion by Theorem~\ref{thm:levychar}.
Finally $T_{\langle M\rangle_t}\ge t$ and $\langle M\rangle$ is constant on $[t,T_{\langle M\rangle_t}]$, so $\beta_{\langle M\rangle_t}=M_{T_{\langle M\rangle_t}}=M_t$.
\end{proof}

\begin{beauty}
Every continuous fair game, however it is built, is Brownian motion watched with a different clock. The variety of continuous martingales is the variety of clocks, and the clock is the quadratic
variation. Doeblin's sealed envelope contained exactly this idea, a quarter of a century before Dambis, Dubins and Schwarz.
\end{beauty}

\begin{remark}
If $\langle M\rangle_\infty<\infty$ with positive probability, the same holds after enlarging the probability space with an independent Brownian motion to run after time $\langle M\rangle_\infty$ (Revuz and Yor, Chapter~V,
Theorem~1.7). In several dimensions, Knight's theorem (1971) says that orthogonal continuous martingales, each with infinite clock, become \emph{independent} Brownian motions after their separate time changes.
\end{remark}

\section{Burkholder's inequality}

How large is a martingale, compared with its clock? For $p=2$ the answer is Doob's inequality and Corollary~\ref{cor:L2bracket}: $\mathbb E\langle M\rangle_t\le\mathbb E\sup_{s\le t}M_s^2\le4\,\mathbb E\langle M\rangle_t$. It\^o's formula
gives the other powers. We first extend Doob's inequality.

\begin{lemma}[Doob's $L^p$ inequality]\label{lem:doobp}
Let $p>1$ and $X$ a continuous non-negative submartingale. Then $\mathbb E\sup_{s\le t}X_s^p\le\bigl(\frac p{p-1}\bigr)^p\mathbb EX_t^p$.
\end{lemma}

\begin{proof}
By dyadic approximation and monotone convergence (as in Theorem~\ref{thm:doobcont}) it suffices to treat a discrete submartingale $X_0,\dots,X_n$, with $X^*=\max_kX_k$ and $\mathbb EX_n^p<\infty$; replacing $X^*$ by $X^*\wedge K$
and letting $K\to\infty$ at the end, we may assume $\mathbb E(X^*)^p<\infty$. The strong form of Volume~4, Theorem~10.15, $\lambda P(X^*\ge\lambda)\le\mathbb E[X_n\mathbf 1_{\{X^*\ge\lambda\}}]$, and Fubini give
\[
\mathbb E(X^*)^p=\int_0^\infty p\lambda^{p-1}P(X^*\ge\lambda)\dd\lambda\le\mathbb E\Bigl[X_n\int_0^{X^*}p\lambda^{p-2}\dd\lambda\Bigr]=\frac p{p-1}\mathbb E\bigl[X_n(X^*)^{p-1}\bigr],
\]
and by H\"older's inequality the right side is at most
\[
\frac p{p-1}(\mathbb EX_n^p)^{1/p}\bigl(\mathbb E(X^*)^p\bigr)^{1-1/p}.
\]
Divide. (With $X^*\wedge K$ in place of $X^*$ the same computation holds, since
$\{X^*\wedge K\ge\lambda\}=\{X^*\ge\lambda\}$ for $\lambda\le K$.)
\end{proof}

\begin{theorem}[Burkholder, 1966]\label{thm:bdg}
Let $p\ge2$ and $C_p=\bigl(\frac{p}{2}(p-1)\bigr)^{p/2}\bigl(\frac p{p-1}\bigr)^{p^2/2}$. For every continuous local martingale $M$ with $M_0=0$ and every $t$,
\[
\mathbb E\sup_{s\le t}|M_s|^p\le C_p\,\mathbb E\langle M\rangle_t^{p/2}.
\]
\end{theorem}

\begin{proof}
Let $M^*_t=\sup_{s\le t}|M_s|$. Stopping at $\inf\{s:|M_s|+\langle M\rangle_s\ge n\}$ and letting $n\to\infty$ at the end (monotone convergence on both sides), we may assume $M$ and $\langle M\rangle$ bounded. Since $x\mapsto|x|^p$
is $C^2$ for $p\ge2$, It\^o's formula gives
\[
|M_t|^p=\int_0^tp|M_s|^{p-1}\operatorname{sgn}(M_s)\dd M_s+\frac{p(p-1)}2\int_0^t|M_s|^{p-2}\dd\langle M\rangle_s .
\]
The stochastic integral has a bounded integrand and a bounded clock, so it is a martingale with mean $0$. Hence, by H\"older's inequality with exponents $\frac p{p-2}$ and $\frac p2$,
\[
\mathbb E|M_t|^p\le\frac{p(p-1)}2\,\mathbb E\bigl[(M^*_t)^{p-2}\langle M\rangle_t\bigr]\le\frac{p(p-1)}2\bigl(\mathbb E(M^*_t)^p\bigr)^{1-2/p}\bigl(\mathbb E\langle M\rangle_t^{p/2}\bigr)^{2/p}.
\]
By Lemma~\ref{lem:doobp} applied to $|M|$, $\mathbb E(M^*_t)^p\le(\frac p{p-1})^p\mathbb E|M_t|^p$. Combining and dividing by $(\mathbb E(M^*_t)^p)^{1-2/p}$ (if it is $0$ there is nothing to prove) gives
$(\mathbb E(M^*_t)^p)^{2/p}\le(\frac p{p-1})^p\frac{p(p-1)}2(\mathbb E\langle M\rangle_t^{p/2})^{2/p}$; raise to the power $\frac p2$.
\end{proof}

\begin{remark}
The reverse inequality $\mathbb E\langle M\rangle_t^{p/2}\le c_p\,\mathbb E\sup_{s\le t}|M_s|^p$ also holds, and both hold for every $p>0$: these are the Burkholder--Davis--Gundy inequalities (1972), proved for
continuous martingales by a ``good $\lambda$'' argument (Revuz and Yor, Chapter~IV, Theorem~4.1). The constant above is far from optimal; its order of growth $p^{p/2}$ is the right one, as Brownian motion shows.
\end{remark}

\section{Tanaka's formula and the local time}

The function $|x|$ is not $C^2$: its second derivative is twice the Dirac mass at $0$. It\^o's formula for $|B_t|$ should therefore contain a term that grows only when $B$ is at $0$. That term is the
\emph{local time}, introduced by L\'evy and made into a process by Hiroshi Tanaka in 1963.

\begin{theorem}[Tanaka's formula]\label{thm:tanaka}
Let $B$ be a Brownian motion. There is a continuous increasing process $L$ with $L_0=0$, the \emph{local time at $0$}, which increases only on the zero set $\mathcal Z=\{t:B_t=0\}$ (it is constant on every interval where
$B$ does not vanish), such that
\[
|B_t|=\int_0^t\operatorname{sgn}(B_s)\dd B_s+L_t\qquad(t\ge0),
\]
where $\operatorname{sgn}(0)=0$. Moreover $\beta_t=\int_0^t\operatorname{sgn}(B_s)\dd B_s$ is a Brownian motion.
\end{theorem}

\begin{proof}
Let $\varphi$ be a smooth even probability density supported in $[-1,1]$, $\varphi_\varepsilon(x)=\varphi(x/\varepsilon)/\varepsilon$, and $f_\varepsilon(x)=\int|x-y|\varphi_\varepsilon(y)\dd y$. Then $f_\varepsilon'(x)=\int\operatorname{sgn}(x-y)\varphi_\varepsilon(y)\dd y=2\Phi_\varepsilon(x)-1$,
with $\Phi_\varepsilon$ the distribution function of $\varphi_\varepsilon$, so $f_\varepsilon''=2\varphi_\varepsilon$ and $f_\varepsilon$ is smooth. Moreover $|f_\varepsilon(x)-|x||\le\varepsilon$, $|f_\varepsilon'|\le1$, and $f_\varepsilon'(x)=\operatorname{sgn}x$ for $|x|\ge\varepsilon$. By
It\^o's formula,
\[
f_\varepsilon(B_t)=f_\varepsilon(0)+\int_0^tf_\varepsilon'(B_s)\dd B_s+\int_0^t\varphi_\varepsilon(B_s)\dd s .
\]
Fix $T$. By Doob's inequality and the isometry,
\[
\mathbb E\sup_{t\le T}\Bigl(\int_0^t\bigl(f_\varepsilon'(B_s)-\operatorname{sgn}B_s\bigr)\dd B_s\Bigr)^2\le4\,\mathbb E\int_0^T\mathbf 1_{\{0<|B_s|<\varepsilon\}}\dd s\xrightarrow[\varepsilon\to0]{}0 .
\]
So $L^\varepsilon_t=\int_0^t\varphi_\varepsilon(B_s)\dd s$ converges, uniformly on $[0,T]$ in $L^2$, to $L_t=|B_t|-\beta_t$; along a sequence $\varepsilon_n\to0$ the convergence is uniform on $[0,T]$ almost surely. Hence $L$ is continuous,
$L_0=0$, and increasing, as a limit of increasing functions. If $B\ne0$ on $[u,v]\subseteq[0,T]$, then $|B|\ge\delta>0$ there, $L^{\varepsilon_n}_v-L^{\varepsilon_n}_u=0$ as soon as $\varepsilon_n<\delta$, and $L_v=L_u$. Let $T\to\infty$
along the integers.

Finally $\beta$ is a continuous local martingale with $\langle\beta\rangle_t=\int_0^t\operatorname{sgn}(B_s)^2\dd s=\int_0^t\mathbf 1_{\{B_s\ne0\}}\dd s=t$, because the zero set has Lebesgue measure zero (Chapter~2). By
Theorem~\ref{thm:levychar}, $\beta$ is a Brownian motion.
\end{proof}

The proof shows that $L_t=\lim_{\varepsilon\to0}\int_0^t\varphi_\varepsilon(B_s)\dd s$ for every smoothing kernel $\varphi$: $L$ is the density at $0$ of the time spent by $B$ near $0$. With $\varphi=\frac12\mathbf 1_{[-1,1]}$ this is L\'evy's
original definition, $L_t=\lim_{\varepsilon\to0}\frac1{2\varepsilon}\lambda\{s\le t:|B_s|<\varepsilon\}$ (Exercise~\ref{ex:occupation}). The next lemma, due to Anatoliy Skorokhod (1961), identifies $L$ from the decomposition alone.

\begin{lemma}[Skorokhod's reflection lemma]\label{lem:skorokhod}
Let $y$ be a continuous function on $[0,\infty[$ with $y_0=0$. Suppose $x=y+\ell$, where $x\ge0$ and $\ell$ is continuous and increasing with $\ell_0=0$, constant on every interval where $x>0$. Then
$\ell_t=\max_{s\le t}(-y_s)$ for all $t$.
\end{lemma}

\begin{proof}
For $s\le t$, $\ell_t\ge\ell_s=x_s-y_s\ge-y_s$, so $\ell_t\ge m_t:=\max_{s\le t}(-y_s)$ (which is $\ge0$ since $y_0=0$). Conversely, if $x$ has no zero in $[0,t]$, then $\ell$ is constant on $[0,t]$ and $\ell_t=0\le m_t$. Otherwise let
$g=\max\{s\le t:x_s=0\}$. Then $\ell$ is constant on $]g,t]$, so by continuity $\ell_t=\ell_g=x_g-y_g=-y_g\le m_t$.
\end{proof}

\begin{theorem}[L\'evy, 1948]\label{thm:levyMB}
Let $B$ be a Brownian motion, $M_t=\max_{s\le t}B_s$, and $L$ the local time at $0$. Then the processes $(M_t-B_t,M_t)_{t\ge0}$ and $(|B_t|,L_t)_{t\ge0}$ have the same law. In particular $M-B$ and $|B|$ have the same
law as processes, and $L_t$ has the law of $|B_t|$ for each $t$.
\end{theorem}

\begin{proof}
Apply Lemma~\ref{lem:skorokhod} to Tanaka's formula, path by path, with $x=|B|$, $y=\beta$ and $\ell=L$: $L_t=\max_{s\le t}(-\beta_s)$. So with $W=-\beta$, a Brownian motion, and $S_t=\max_{s\le t}W_s$,
\[
\bigl(|B_t|,L_t\bigr)=\bigl(S_t-W_t,S_t\bigr)\qquad\text{for all }t\text{, almost surely.}
\]
The right side is the same measurable function of the Brownian motion $W$ as $(M-B,M)$ is of $B$, so the two pairs have the same law.
\end{proof}

\begin{corollary}[the dimension of the zero set]\label{cor:zerodim}
Almost surely, $\dim_H(\mathcal Z\cap[0,1])=\frac12$.
\end{corollary}

\begin{proof}
In the proof of Theorem~\ref{thm:levyMB}, $B_t=0$ if and only if $W_t=S_t$: the zero set of $B$ is \emph{equal}, path by path, to the record set of the Brownian motion $W$. By Theorem~\ref{thm:dimension} applied to $W$,
its intersection with $[0,1]$ has dimension $\frac12$ almost surely.
\end{proof}

\begin{figure}[ht]
\centering
\fig{ch04-tanaka}
\caption{Top: a Brownian path $B$ on $[0,1]$. Bottom: $|B|$ (thin), its local time at $0$ (thick), and the running maximum of $W=-\int\operatorname{sgn}(B)\dd B$ (dashed, visible only where it differs from the
local time by the discretisation), and L\'evy's occupation estimate $\frac1{2\varepsilon}\lambda\{s\le t:|B_s|<\varepsilon\}$ with $\varepsilon=0.01$ (dotted). The local time grows only on the zero set, and $|B|=S-W$: the reflected path is a path seen from its running maximum.}
\label{fig:tanaka}
\end{figure}

So the zero set carries a natural measure, $\mathrm dL$, which plays for $\mathcal Z$ the role that $\mathrm dM$ played for the record set in Chapter~2.
In Chapter~7 the inverse local time $\tau_\ell=\inf\{t:L_t>\ell\}$ will reappear as the $\frac12$-stable subordinator, the clock that counts excursions.

\section{Exercises}

\begin{exercise}[$\star$]\label{ex:bdb}
(a) Show that $\int_0^tB_s\dd B_s=\frac12(B_t^2-t)$. (b) Show that the right-point sums $\sum B_{t_i}\Delta B_i$ converge in probability to $\frac12(B_t^2+t)$ and the midpoint sums $\sum\frac12(B_{t_{i-1}}+B_{t_i})\Delta B_i$ to
$\frac12B_t^2$. Which of the three limits is a martingale?
\end{exercise}

\begin{exercise}[$\star$]\label{ex:moments}
(a) With It\^o's formula, show that $\mathbb EB_t^4=3t^2$ and, more generally, $\mathbb EB_t^{2k}=k(2k-1)\int_0^t\mathbb EB_s^{2k-2}\dd s$. (b) Let $N_t=B_t^2-t$. Show that $\langle N\rangle_t=4\int_0^tB_s^2\dd s$ and recover $\operatorname{Var}B_t^2=2t^2$.
\end{exercise}

\begin{exercise}[$\star$]\label{ex:ibp}
(a) Prove the integration by parts formula $X_tY_t=X_0Y_0+\int_0^tX\dd Y+\int_0^tY\dd X+\langle X,Y\rangle_t$ for continuous semimartingales. (b) Deduce $tB_t=\int_0^ts\dd B_s+\int_0^tB_s\dd s$, and that
$\int_0^1B_s\dd s$ is $N(0,\frac13)$.
\end{exercise}

\begin{exercise}[$\star\star$]\label{ex:wienerito}
Let $f\in L^2[0,\infty[$ be deterministic and $X_t=\int_0^tf(s)\dd B_s$ (It\^o's integral). (a) Show that $X_t$ equals the Wiener integral $\int f\mathbf 1_{[0,t]}\dd B$ of Theorem~\ref{thm:wienerint}, so that $X$ is a continuous
Gaussian martingale with independent increments. (b) Show that $\exp\bigl(i\xi X_t+\frac{\xi^2}2\int_0^tf^2\bigr)$ is a martingale, and recover the law of $X_t$. (c) If $f(s)=e^{s}$, express $X$ as a time-changed Brownian motion.
\end{exercise}

\begin{exercise}[$\star\star$]\label{ex:rotation}
(a) Let $B$ be a $d$-dimensional Brownian motion and $O$ a progressive process with values in the orthogonal matrices. Show that $W_t=\int_0^tO_s\dd B_s$ is a $d$-dimensional Brownian motion. (b) Let
$\beta_t=\int_0^t\frac{B_s}{|B_s|}\cdot\dd B_s$ (with $\frac00=0$). Show that $\beta$ is a one-dimensional Brownian motion and that $|B_t|^2=2\int_0^t|B_s|\dd\beta_s+dt$. (The process $|B|^2$ is the
\emph{squared Bessel process} of dimension $d$; it reappears in Chapter~6.)
\end{exercise}

\begin{exercise}[$\star\star$]\label{ex:occupation}
Show that $\frac1{2\varepsilon}\lambda\{s\le t:|B_s|<\varepsilon\}\to L_t$ in probability as $\varepsilon\to0$. More generally, show that $\int_0^tg_\varepsilon(B_s)\dd s\to L_t$ for every probability density $g$ supported in $[-1,1]$, smooth or
not, where $g_\varepsilon(x)=g(x/\varepsilon)/\varepsilon$.
\end{exercise}

\begin{exercise}[$\star$]\label{ex:locallaw}
(a) Show that $\mathbb EL_t=\mathbb E|B_t|=\sqrt{2t/\pi}$, in two ways: from Tanaka's formula and from Theorem~\ref{thm:levyMB}. (b) Show that $L_t\to\infty$ almost surely as $t\to\infty$. (c) Show that $P(L_t=0)=0$ for $t>0$,
and interpret.
\end{exercise}

\begin{exercise}[$\star\star\star$]\label{ex:martslln}
Let $M$ be a continuous local martingale with $M_0=0$. (a) Show that $M_t$ converges almost surely, as $t\to\infty$, on the event $\{\langle M\rangle_\infty<\infty\}$. (b) If $\langle M\rangle_\infty=\infty$ almost surely, show that
$M_t/\langle M\rangle_t\to0$ almost surely, and that $\limsup_t M_t=+\infty$ and $\liminf_tM_t=-\infty$. (c) Apply (b) to $M_t=\int_0^tB_s\dd B_s$ to show that $\int_0^tB_s^2\dd s\to\infty$, a fact that is easy to see
directly, and that $B_t^2-t$ changes sign infinitely often.
\end{exercise}

\section{Hints}

\begin{hint}{ex:bdb}
(a) $f(x)=x^2$. (b) The right sum minus the left sum is $\sum(\Delta B_i)^2$.
\end{hint}

\begin{hint}{ex:moments}
(a) The stochastic integral $\int_0^t2kB_s^{2k-1}\dd B_s$ has an integrand in $L^2$ on $[0,t]$, so it has mean zero. (b) $N=2\int B\dd B$; use Proposition~\ref{prop:itoprops}(a).
\end{hint}

\begin{hint}{ex:ibp}
(a) It\^o's formula for $f(x,y)=xy$, or polarisation from $X^2$. (b) $X_t=t$ has no bracket. The integral $\int_0^ts\dd B_s$ is a Wiener integral (Exercise~\ref{ex:wienerito}).
\end{hint}

\begin{hint}{ex:wienerito}
(a) Both integrals are linear isometries from $L^2[0,t]$ that agree on step functions. (c) $\langle X\rangle_t=\frac12(e^{2t}-1)$.
\end{hint}

\begin{hint}{ex:rotation}
(a) Compute $\langle W^k,W^l\rangle$ with Proposition~\ref{prop:itoprops}(a) and apply L\'evy's characterisation. (b) $\langle\beta\rangle_t=\int_0^t\mathbf 1_{\{B_s\ne0\}}\dd s$; It\^o's formula for $|x|^2$.
\end{hint}

\begin{hint}{ex:occupation}
Squeeze $g$ between smooth kernels: for $\delta>0$ there are smooth $\varphi^-\le g\le\varphi^+$, supported in $[-1-\delta,1+\delta]$, with $\int\varphi^\pm$ within $\delta$ of $1$ (for $g=\frac12\mathbf 1_{[-1,1]}$ this is easy to draw), and the proof of
Theorem~\ref{thm:tanaka} works for any non-negative smooth kernel with compact support, giving the limit $(\int\varphi)L_t$.
\end{hint}

\begin{hint}{ex:locallaw}
(b) In the proof of Theorem~\ref{thm:levyMB}, $L$ is the running maximum of a Brownian motion $W$. (c) Use the equality in law of (a).
\end{hint}

\begin{hint}{ex:martslln}
(a) Stop at $\sigma_n=\inf\{t:\langle M\rangle_t\ge n\}$. (b) Dambis--Dubins--Schwarz and $\beta_s/s\to0$, and the oscillation of Brownian motion at infinity (Corollary~\ref{cor:infinity}).
\end{hint}

\section{Solutions}

\begin{solution}{ex:bdb}
\idea The three sums differ only in where the integrand is evaluated, and the differences are sums of squared increments, which do not vanish: they converge to the clock. Only the left point makes each term a
fair bet, because the integrand is then known before the increment.
\writeup (a) It\^o's formula with $f(x)=x^2$: $B_t^2=2\int_0^tB\dd B+t$. (b) $\sum B_{t_i}\Delta B_i=\sum B_{t_{i-1}}\Delta B_i+\sum(\Delta B_i)^2\to\frac12(B_t^2-t)+t$ by Proposition~\ref{prop:riemann}, and the midpoint sum is the
average of the left and right sums, $\frac12B_t^2$; alternatively it telescopes exactly: $\sum\frac12(B_{t_{i-1}}+B_{t_i})(B_{t_i}-B_{t_{i-1}})=\frac12\sum(B_{t_i}^2-B_{t_{i-1}}^2)=\frac12B_t^2$. Only the It\^o limit $\frac12(B_t^2-t)$ is a
martingale. The midpoint convention, \emph{Stratonovich's}, obeys the ordinary chain rule, at the price of losing the martingale property.
\end{solution}

\begin{solution}{ex:moments}
\idea It\^o's formula turns a moment into the integral of a lower moment, because the stochastic integral has mean zero. The bracket of an integral is the integral of the squared integrand against the clock.
\writeup (a) $B_t^{2k}=\int_0^t2kB_s^{2k-1}\dd B_s+k(2k-1)\int_0^tB_s^{2k-2}\dd s$. The integrand of the stochastic integral satisfies $\mathbb E\int_0^tB_s^{4k-2}\dd s<\infty$, so by Theorem~\ref{thm:itoloc} the integral is in $\mathbb H^2$ on
$[0,t]$ and has mean $0$. Taking expectations, $\mathbb EB_t^{2k}=k(2k-1)\int_0^t\mathbb EB_s^{2k-2}\dd s$; for $k=2$, $6\int_0^ts\dd s=3t^2$. (b) $N_t=2\int_0^tB\dd B$ by Exercise~\ref{ex:bdb}, so $\langle N\rangle_t=4\int_0^tB_s^2\dd s$, and
$\operatorname{Var}B_t^2=\mathbb EN_t^2=\mathbb E\langle N\rangle_t=4\int_0^ts\dd s=2t^2$, in agreement with $3t^2-t^2$.
\end{solution}

\begin{solution}{ex:ibp}
\idea The product is a smooth function of two semimartingales with constant second derivatives; the cross term of It\^o's formula is the bracket.
\writeup (a) For $f(x,y)=xy$, $\partial_x\partial_yf=1$ and $\partial_x^2f=\partial_y^2f=0$, so It\^o's formula gives $\mathrm d(XY)=X\dd Y+Y\dd X+\frac12\cdot2\dd\langle X,Y\rangle$. (b) With $X_t=t$ and $Y=B$, $\langle X,B\rangle=0$:
$tB_t=\int_0^ts\dd B_s+\int_0^tB_s\dd s$. At $t=1$, $\int_0^1B_s\dd s=\int_0^1(1-s)\dd B_s$, a Wiener integral, Gaussian with variance $\int_0^1(1-s)^2\dd s=\frac13$.
\end{solution}

\begin{solution}{ex:wienerito}
\idea For a deterministic integrand, It\^o's construction and Wiener's coincide: both start from step functions and extend by isometry. Then L\'evy's computation with $\exp(i\xi X)$ shows the Gaussian law directly, and
Dambis--Dubins--Schwarz turns the process into Brownian motion with a deterministic clock.
\writeup (a) For a step function $f=\sum c_i\mathbf 1_{]t_i,t_{i+1}]}$, both integrals equal $\sum c_i(B_{t_{i+1}}-B_{t_i})$. Both are isometries from $L^2[0,t]$ into $L^2(P)$ (Theorem~\ref{thm:wienerint} and
Theorem~\ref{thm:itoloc}), and step functions are dense in $L^2[0,t]$, so they agree. So $X$ is a continuous martingale, and Gaussian with independent increments by Theorem~\ref{thm:wienerint} (indicators of disjoint
intervals times $f$ are orthogonal). (b) As in the proof of Theorem~\ref{thm:levychar}, It\^o's formula for $\exp(i\xi x+\frac{\xi^2}2y)$ with $X$ and $y=\langle X\rangle_t=\int_0^tf^2$ gives a bounded local martingale on $[0,t]$, hence
$\mathbb Ee^{i\xi X_t}=e^{-\frac{\xi^2}2\int_0^tf^2}$: $X_t$ is $N(0,\int_0^tf^2)$. (c) $\langle X\rangle_t=\int_0^te^{2s}\dd s=\frac12(e^{2t}-1)$, which tends to infinity, so by Theorem~\ref{thm:dds} $X_t=\beta_{(e^{2t}-1)/2}$ for a Brownian
motion $\beta$. (Then $e^{-t}X_t$ is an Ornstein--Uhlenbeck process; compare Exercise~\ref{ex:ou}.)
\end{solution}

\begin{solution}{ex:rotation}
\idea L\'evy's characterisation only asks for the clocks, and a rotation does not change clocks: the brackets of the rotated motion are $OO^{\mathsf T}=I$ times $t$.
\writeup (a) $W^k=\sum_m\int O^{km}\dd B^m$, so by Proposition~\ref{prop:itoprops}(a) and $\langle B^m,B^n\rangle_t=\delta_{mn}t$, $\langle W^k,W^l\rangle_t=\int_0^t\sum_mO^{km}_sO^{lm}_s\dd s=\int_0^t(O_sO_s^{\mathsf T})_{kl}\dd s=\delta_{kl}t$.
By Theorem~\ref{thm:levychar}, $W$ is a Brownian motion. (b) $\beta=\sum_k\int\frac{B^k}{|B|}\dd B^k$ is a continuous local martingale with $\langle\beta\rangle_t=\int_0^t\sum_k\frac{(B^k_s)^2}{|B_s|^2}\mathbf 1_{\{B_s\ne0\}}\dd s=t$, since
$\{s:B_s=0\}\subseteq\{s:B^1_s=0\}$ has measure zero. So $\beta$ is a Brownian motion. It\^o's formula for $|x|^2$ gives $|B_t|^2=\sum_k2\int_0^tB^k\dd B^k+dt$, and $\sum_kB^k\dd B^k=|B|\dd\beta$ by
Proposition~\ref{prop:itoprops}(b) (on $\{B=0\}$ both integrands vanish).
\end{solution}

\begin{solution}{ex:occupation}
\idea The local time does not care how we smooth: any kernel of total mass one sees the same density at $0$. Sandwiching a rough kernel between smooth ones transfers the result.
\writeup The proof of Theorem~\ref{thm:tanaka} used only that $\varphi$ is smooth, non-negative, supported in a compact interval $[-a,a]$, with $\int\varphi=1$ (and then $|f_\varepsilon(x)-|x||\le a\varepsilon$). For a smooth non-negative $\psi$
with compact support and mass $m>0$, apply it to $\psi/m$: $\int_0^t\psi_\varepsilon(B_s)\dd s\to mL_t$, uniformly on compacts in probability. Given $g$ and $\delta$, choose smooth $\varphi^-\le g\le\varphi^+$ with masses $m^\pm$ in
$[1-\delta,1+\delta]$. Then $\int_0^t\varphi^-_\varepsilon(B)\le\int_0^tg_\varepsilon(B)\le\int_0^t\varphi^+_\varepsilon(B)$, and the outer terms tend to $m^\pm L_t$. So, for $\eta>0$, with probability tending to $1$,
$(1-\delta)L_t-\eta\le\int_0^tg_\varepsilon(B_s)\dd s\le(1+\delta)L_t+\eta$. Since $\delta$ and $\eta$ are arbitrary and $L_t$ is finite, $\int_0^tg_\varepsilon(B)\dd s\to L_t$ in probability. With $g=\frac12\mathbf 1_{[-1,1]}$,
$\int_0^tg_\varepsilon(B_s)\dd s=\frac1{2\varepsilon}\lambda\{s\le t:|B_s|<\varepsilon\}$, up to the null set $\{|B_s|=\varepsilon\}$.
\end{solution}

\begin{solution}{ex:locallaw}
\idea The local time is the compensator of $|B|$, so they have the same mean; L\'evy's theorem says more, that they have the same law at each time.
\writeup (a) $\beta$ is a martingale with mean $0$ (it is in $\mathbb H^2$ on $[0,t]$), so $\mathbb EL_t=\mathbb E|B_t|=\sqrt t\,\mathbb E|Z|=\sqrt{2t/\pi}$. By Theorem~\ref{thm:levyMB}, $L_t$ has the law of $|B_t|$ directly. (b) $L_t=S_t$, the
running maximum of a Brownian motion $W$, which tends to infinity because $\limsup W_t=+\infty$ (Corollary~\ref{cor:infinity}). (c) $P(L_t=0)=P(|B_t|=0)=0$. The local time starts growing immediately: $B$ returns to $0$
infinitely often in every interval $[0,\varepsilon]$ (Corollary~\ref{cor:oscillate}), and these returns accumulate positive local time, as they must, since $W$ has records immediately after time $0$.
\end{solution}

\begin{solution}{ex:martslln}
\idea On a finite clock, the martingale is a Brownian motion observed on a bounded time interval, which converges at the end. On an infinite clock, it is a whole Brownian path, and the law of large numbers for
Brownian motion transfers through the time change.
\writeup (a) $M^{\sigma_n}$ has bracket at most $n$, so it is bounded in $L^2$ and converges almost surely (Corollary~\ref{cor:L2bracket}, Theorem~\ref{thm:L2cont}). On $\{\langle M\rangle_\infty<n\}$, $\sigma_n=\infty$ and $M=M^{\sigma_n}$.
Taking the union over $n$ gives (a). (b) By Theorem~\ref{thm:dds}, $M_t=\beta_{\langle M\rangle_t}$ with $\beta$ a Brownian motion, and $\langle M\rangle_t\to\infty$. Since $\beta_s/s\to0$ (Volume~4, Proposition~11.7), $M_t/\langle M\rangle_t\to0$; and
since $\limsup_s\beta_s=+\infty$, $\liminf_s\beta_s=-\infty$ (Corollary~\ref{cor:infinity}) and $\langle M\rangle_t$ runs through all of $[0,\infty[$ continuously, $M$ has the same $\limsup$ and $\liminf$. (c) $M=\int B\dd B$ has
$\langle M\rangle_t=\int_0^tB_s^2\dd s$. If this stayed bounded with positive probability, $M_t=\frac12(B_t^2-t)$ would converge there by (a), which is impossible since $B_t^2-t\to-\infty$ along times when $B$ is at $0$; so
$\langle M\rangle_\infty=\infty$ almost surely (that is the direct fact), and then by (b) $\limsup M_t=+\infty$ and $\liminf M_t=-\infty$: $B_t^2-t$ changes sign infinitely often.
\end{solution}

\section*{Sources and further reading}

\begin{itemize}
\item \emph{History.} K.~It\^o, ``Stochastic integral'', \emph{Proceedings of the Imperial Academy Tokyo} 20 (1944), 519--524, and ``On a formula concerning stochastic differentials'', \emph{Nagoya Mathematical Journal} 3
(1951), 55--65. K.~It\^o, \emph{Selected Papers}, edited by D.~W.~Stroock and S.~R.~S.~Varadhan, Springer, 1987, with It\^o's foreword. W.~Doeblin, ``Sur l'\'equation de Kolmogoroff'', pli cachet\'e of 1940,
\emph{Comptes Rendus de l'Acad\'emie des Sciences}, S\'erie~I, 331 (2000), special issue. B.~Bru and M.~Yor, ``Comments on the life and mathematical legacy of Wolfgang Doeblin'', \emph{Finance and Stochastics} 6
(2002), 3--47. P.~L\'evy, \emph{Processus stochastiques et mouvement brownien}, Gauthier-Villars, 1948. H.~Tanaka, ``Note on continuous additive functionals of the 1-dimensional Brownian path'',
\emph{Zeitschrift f\"ur Wahrscheinlichkeitstheorie} 1 (1963), 251--257. A.~V.~Skorokhod, ``Stochastic equations for diffusion processes in a bounded region'', \emph{Theory of Probability and its Applications} 6 (1961),
264--274. K.~E.~Dambis, ``On the decomposition of continuous submartingales'', \emph{Theory of Probability and its Applications} 10 (1965), 401--410; L.~E.~Dubins and G.~Schwarz, ``On continuous martingales'',
\emph{Proceedings of the National Academy of Sciences USA} 53 (1965), 913--916. D.~L.~Burkholder, ``Martingale transforms'', \emph{Annals of Mathematical Statistics} 37 (1966), 1494--1504;
D.~L.~Burkholder, B.~J.~Davis and R.~F.~Gundy, ``Integral inequalities for convex functions of operators on martingales'', \emph{Proceedings of the Sixth Berkeley Symposium} 2 (1972), 223--240.
\item J.-F.~Le~Gall, \emph{Brownian Motion, Martingales, and Stochastic Calculus}, Springer, 2016, Chapters~5 and~9. The construction followed here, through the characterisation by brackets.
\item D.~Revuz and M.~Yor, \emph{Continuous Martingales and Brownian Motion}, 3rd ed., Springer, 1999, Chapters~IV, V and VI. Time changes, the BDG inequalities, local times in full.
\item I.~Karatzas and S.~E.~Shreve, \emph{Brownian Motion and Stochastic Calculus}, 2nd ed., Springer, 1991, Chapters~3 and~6.
\item P.~M\"orters and Y.~Peres, \emph{Brownian Motion}, Cambridge University Press, 2010, Chapters~6 and~7. Local times, L\'evy's theorem and the zero set from the point of view of path properties.
\end{itemize}
